\documentclass[11pt]{article}

\usepackage[T1]{fontenc}
\usepackage{lmodern}
\usepackage{amsmath,amssymb,amsthm}
\usepackage[margin=1in]{geometry}
\usepackage{booktabs,longtable}
\usepackage{enumitem}
\usepackage{xcolor}
\usepackage{microtype}
\usepackage{seqsplit}
\usepackage[hidelinks]{hyperref}

\hypersetup{
  pdftitle={Detailed errata and statement audit for the chemotaxis--logistic trilogy},
  pdfauthor={ShenWork formalization project}
}

\newcommand{\R}{\mathbb{R}}
\newcommand{\Om}{\Omega}
\newcommand{\Tmax}{T_{\max}}
\newcommand{\norm}[1]{\left\lVert #1\right\rVert}
\newcommand{\abs}[1]{\left\lvert #1\right\rvert}
\newcommand{\code}[1]{\nolinkurl{#1}}
\newcommand{\False}{\textsc{False statement}}
\newcommand{\Gap}{\textsc{Proof gap}}
\newcommand{\Clarify}{\textsc{Clarification}}
\newcommand{\Interface}{\textsc{Interface only}}
\setlength{\emergencystretch}{3em}
\sloppy

\newtheorem{proposition}{Proposition}
\newtheorem{remark}{Remark}

\title{Detailed Errata and Statement Audit\\
\large Three chemotaxis--logistic papers formalized in Lean~4}
\author{\code{ShenWork} formalization project}
\date{27 July 2026}

\begin{document}
\maketitle

\begin{abstract}
This document records statement defects and proof gaps found while formalizing
three papers on parabolic--elliptic chemotaxis systems.  It deliberately
separates four logically different outcomes: an explicit counterexample to a
printed conclusion, a gap in the printed proof without a counterexample to the
conclusion, a wording clarification, and a defect in an earlier Lean interface
that is not an error in the paper.  For every mathematical counterexample we
give the calculation, the corrected statement, and the exact Lean certificate.
The repository is pinned to Lean~4.30.0 and Mathlib~v4.30.0.  The cited
certificates are \code{sorry}-free and their headline axiom audits contain only
\code{propext}, \code{Classical.choice}, and \code{Quot.sound}.
\end{abstract}

\tableofcontents

\section{Scope, sources, and classification}

The audited sources are:
\begin{enumerate}[leftmargin=2em]
\item Wenxian Shen, \emph{Existence, uniqueness, stability, and monotonicity
  of traveling waves for repulsion/attraction chemotaxis models with logistic
  type source}, arXiv:2605.04401v1 (the ``traveling-wave paper'');
\item Le Chen, Ian Ruau, and Wenxian Shen,
  \emph{Chemotaxis models with signal-dependent sensitivity and a
  logistic-type source, I: Boundedness and global existence},
  arXiv:2512.14858v1 (``Part~I'');
\item Le Chen, Ian Ruau, and Wenxian Shen,
  \emph{Chemotaxis models with signal-dependent sensitivity and a
  logistic-type source, II: Persistence and stabilization},
  arXiv:2604.02599v1 (``Part~II'').
\end{enumerate}

\noindent The labels used below have precise meanings.
\begin{description}[leftmargin=2.1em,style=nextline]
\item[\False] The printed hypotheses admit an explicit object for which the
  printed conclusion fails.
\item[\Gap] A displayed step is invalid or a required hypothesis is missing
  from the proof.  No counterexample to the theorem itself is asserted.
\item[\Clarify] The mathematical intent is recoverable, but a literal reading
  is stronger or less precise than the continuation or regularity framework
  supports.
\item[\Interface] An earlier Lean definition quantified over more objects, or
  stored less semantic information, than the paper.  A refutation of that Lean
  definition is not evidence against the paper.
\end{description}

\begin{center}
\renewcommand{\arraystretch}{1.18}
\begin{longtable}{@{}p{0.06\textwidth}p{0.17\textwidth}p{0.44\textwidth}p{0.22\textwidth}@{}}
\toprule
\# & Source & Finding & Status\\
\midrule
\endhead
E1 & Traveling wave & Theorems 1.2--1.3: seven defects in the Section 5
  weighted-energy argument & \Gap; corrected proof window\\
E2 & Traveling wave & Lemma 4.2: the constant subsolution is not valid for
  every member of the printed trap set & \False\\
E3 & Part I & Theorem 1.2 admits \(a>0,b=0\), where constant solutions grow
  exponentially & \False\\
E4 & Part I & Theorem 1.3(iv) invokes an elliptic estimate outside its exponent
  domain & \Gap\\
E5 & Part I & Proposition 1.1: the endpoint alternative belongs to the maximal
  continuation, not an arbitrary finite witness & \Clarify\\
E6 & Part I & Lemma 2.7 controls only \(t>0\) but concludes a bound including
  \(t=0\) & \False\ as written\\
E7 & Part II & Theorem 2.1(1) admits the pure-decay regime \(a=0<b\) & \False\\
E8 & Part II & Theorems 2.2--2.5 reuse an all-time \(C^1\) estimate from merely
  continuous or \(L^\infty\)-small data & \Clarify; ill-posed at \(t=0\)\\
\bottomrule
\end{longtable}
\end{center}

\section{Traveling-wave paper}

\subsection{E1: Theorems 1.2--1.3, weighted stability}

\paragraph{Printed claim.}
Let
\[
  \kappa=\frac{c-\sqrt{c^2-4}}2,\qquad
  \kappa^2-c\kappa+1=0.
\]
Theorem~1.2 states weighted stability for
\[
  \kappa<\eta<\frac{1}{1+\abs{\chi}^{1/6}},
\]
and (1.21) is printed with the laboratory-coordinate weight
\(e^{2\eta x}\).  Theorem~1.3 uses the same stability theorem in its uniqueness
argument.

\paragraph{The load-bearing root error.}
The coefficient produced by (5.31) is
\[
  q(\eta)=\eta^2-(c-A)\eta+(1+B),
\]
where the displayed quantities in (5.32)--(5.33) have \(A\ge0\) and, for
\(\chi\ne0\), \(B>0\).  Evaluation at the unperturbed tail exponent is exact:
\[
  q(\kappa)
   =(\kappa^2-c\kappa+1)+A\kappa+B
   =A\kappa+B>0.
\]
Therefore \(\kappa\) is not between the two roots of \(q\).  Under the speed
condition it lies below the lower perturbed root
\[
  \kappa^-=\frac{(c-A)-\sqrt{(c-A)^2-4(1+B)}}2.
\]
Continuity of \(q\) then gives \(\eta>\kappa\), arbitrarily close to
\(\kappa\), for which \(q(\eta)>0\).  The global coefficient estimate written
in Section~5 proves decay only for \(\eta>\kappa^-\), not for every
\(\eta>\kappa\).  This is a proof gap, not a constructed unstable solution:
a new localized-coercivity or weighted spectral argument could conceivably
recover the larger interval.

\paragraph{Coordinate and exponential signs.}
The energy actually estimated after changing to \(z=x-ct\) is
\[
 E_{\mathrm{move}}(t)=
   \int_{\R}e^{2\eta z}\abs{u(t,z+ct)-U(z)}^2\,dz.
\]
The printed (1.21) is
\[
 E_{\mathrm{lab}}(t)=
   \int_{\R}e^{2\eta x}\abs{u(t,x)-U(x-ct)}^2\,dx.
\]
Changing variables \(x=z+ct\) yields
\[
  E_{\mathrm{lab}}(t)=e^{2\eta ct}E_{\mathrm{move}}(t).
\]
Thus decay of the moving-frame quantity does not by itself imply the printed
laboratory-weighted limit.  The correct laboratory notation is
\(e^{2\eta(x-ct)}\).  In addition, the paper declares
\(\lambda=q(\eta)<0\) and then uses \(e^{-\lambda t}\) as a decaying factor.
For this sign convention the decaying factor is \(e^{\lambda t}\), or,
equivalently, \(e^{-\rho t}\) with \(\rho=-\lambda>0\).

\paragraph{Four coefficient errors.}
Writing \(w=u-U\), \(z=v-V\), the elliptic relation
\(v_{xx}=v-u^\gamma\) gives the zero-order chemotactic flux difference
\[
  (v a_m-a_{m+\gamma})w+U^m z.
\]
After multiplication by \(-\chi\), the perturbation equation contains
\[
 -\chi(v a_m-a_{m+\gamma}+b_2)w-\chi b_4z.
\]
Equations (5.18)--(5.19) reverse both the \(a_{m+\gamma}\) and \(b_4z\)
signs.  These signs can be repaired by absolute-value estimates.  Three
further estimates must also be corrected:
\begin{align*}
 \abs{\chi b_1 w_xw}
   &\le \frac12w_x^2+\frac12\abs{\chi}^2b_1^2w^2,
   &&\text{not a term linear in \(b_1\)},\\
 \int_\R V^2
   &\le \frac{\gamma^2M^{2(\gamma-1)}}{(1-\eta)^2}\int_\R W^2,\\
 \int_\R V_x^2
   &\le \frac{\gamma^2M^{2(\gamma-1)}}{1-\eta^2}\int_\R W^2.
\end{align*}
The factor \(M^{2(\gamma-1)}\) is absent from (5.29)--(5.30).  Finally,
Lemma~5.2 proves only \(U_x/U\le C\), while (5.23) uses
\(\abs{U_x/U}\le C\).  The latter requires a two-sided logarithmic-derivative
estimate.

\paragraph{Corrected form and Lean evidence.}
The formalization retains the full corrected budget \(A_\chi,B_\chi\), the
root window
\[
  \kappa^-_{c,\chi,m,\alpha,\gamma}<\eta
  <\frac{1}{1+\abs{\chi}^{1/6}},
\]
and moving-frame weighted convergence.  The main certificates are:
\begin{itemize}[leftmargin=2em]
\item \code{paper531_kappa_not_between_perturbed_roots},
  \code{paper531_kappa_lt_rootMinus}, and
  \code{paper531_positive_inside_stated_weight_window} in
  \code{Paper1/Theorem12RootObstruction.lean};
\item \code{laboratoryWeightedL2Energy_eq_exp_mul_coMoving} in
  \code{Paper1/Theorem12CoordinateAudit.lean};
\item the corrected flux identity in
  \code{Paper1/Theorem12MeanCoefficients.lean};
\item the squared Young term and resolver cap in
  \code{Paper1/Theorem12WeightedEnergy.lean};
\item the two-sided Riccati fencing estimate in
  \code{Paper1/Theorem12LogDerivative.lean}.
\end{itemize}
The corrected headline theorems are
\code{paper1_Theorem_1_2_chi_pos_unconditional},
\code{paper1_Theorem_1_2_chi_zero_unconditional}, and
\code{paper1_Theorem_1_2_chi_neg_unconditional}.  Their quantified weight
range begins at the corrected \code{paper531RootMinus}.

\subsection{E2: Lemma 4.2, constant subsolution}

\paragraph{Printed claim.}
For \(0<\kappa<\widetilde\kappa\) in the displayed parameter range,
\(M\ge1\), \(D>D_{M,\kappa,\widetilde\kappa,\chi,m,\gamma}\), every
\(u\in E_{\kappa,M,T}\), and every
\[
  0<d\le d_{\kappa,\widetilde\kappa,D,\chi},
\]
Lemma~4.2(2) asserts that the constant function \(W=d\) is a subsolution on
the whole line.

\paragraph{Counterexample.}
Suppress time and take the admissible trap profile
\[
 u(y)=\min\{1,e^{-y/2}\}\min\{1,\abs{y+1}\}.
\]
It is continuous, nonnegative, bounded by the printed upper barrier for
\(\kappa=\frac12\), \(M=1\), and satisfies
\[
  u(-1)=0,\qquad u(0)=1.
\]
For \(m=\alpha=\gamma=1\), the frozen elliptic response at \(-1\) is
\[
  V_0=\frac12\int_\R e^{-\abs{-1-y}}u(y)\,dy>0.
\]
Choose
\[
  \chi=\frac2{V_0},\qquad
  \kappa=\frac12,\quad\widetilde\kappa=1,\quad M=1,\quad c=\frac52,
\]
and any \(D\) one unit above the printed \(D\)-threshold.  The printed
constant threshold is positive; let \(d\) be half of it.  At \(x=-1\), all
spatial derivatives of \(W=d\) vanish, while the frozen operator equals
\[
\begin{aligned}
 \mathcal A(d;u)(-1)
  &= -\chi d\bigl(V_0-u(-1)\bigr)+d(1-d)\\
  &= -2d+d-d^2\\
  &= -d(1+d)<0.
\end{aligned}
\]
A subsolution requires \(\mathcal A(d;u)\ge0\).  Hence the literal universal
statement is false.  The problem is not the nonsmooth interface of a minimum:
the counterexample uses the separately stated constant subsolution and tests
it at a regular point.

\paragraph{Repair.}
The missing content is a pointwise or uniform control of the frozen elliptic
forcing.  A transparent sufficient condition on a trap bounded by \(M\) is
\[
  \abs{\chi}\,d^{m-1}M^\gamma\le 1-d^\alpha.
\]
For example, the formalization proves the constant branch from
\[
  \abs{\chi}M^\gamma\le\frac12,\qquad 0<d\le\frac12.
\]
It also proves the \(\chi=0\) lower- and constant-subsolution slices under
the corresponding \(D\)-scale hypotheses.  The refutation is
\code{not_Lemma_4_2} in \code{Paper1/Statements.lean}; the strengthened
lemmas in the same file include
\code{constant_subsolution_frozenWaveOperator_nonneg_of_half_bound_trap} and
\code{Lemma_4_2_chi_zero_subsolutions_of_D_ge_one}.

\section{Part I: boundedness and global existence}

\subsection{E3: Theorem 1.2, missing reaction guard}

\paragraph{Printed range.}
The theorem assumes \(a,b\ge0\) and \(\beta\ge1\), which includes
\(a>0,b=0\), and asserts uniform boundedness in both its \(0<m<1\) and
small-sensitivity \(m=1\) branches.

\paragraph{Counterexample.}
On any smooth bounded Neumann domain, let \(c_0>0\) and set
\[
  u(t,x)=c_0e^{at},\qquad
  v(t,x)=\frac{\nu}{\mu}u(t,x)^\gamma.
\]
The functions are spatially constant, so diffusion, chemotactic divergence,
and all boundary fluxes vanish.  The elliptic equation holds, and the
parabolic equation reduces to \(u_t=au\).  Thus this is a positive global
classical solution for every \(m,\beta,\chi_0\), but
\[
  \norm{u(t)}_\infty=c_0e^{at}\longrightarrow\infty.
\]
Equivalently, every positive solution in this branch satisfies
\[
  \frac{d}{dt}\int_\Om u=a\int_\Om u.
\]
The counterexample already satisfies the \(m=1,\beta=1,\chi_0=0\) smallness
condition.

\paragraph{Correction and certificate.}
Under the standing nonnegativity assumptions, the maximal natural correction
is
\[
  a=0\quad\text{or}\quad b>0,
\]
which excludes exactly \(a>0,b=0\).  The more conservative split used
verbatim in the paper's proof is
\((a=b=0)\lor(a>0\land b>0)\).  The Lean refutations are
\code{not_Theorem_1_2_intervalDomain_of_a_pos_b_zero} and
\code{not_Theorem_1_2_intervalDomain_when_a_pos_b_zero} in
\code{Paper2/IntervalDomainTheorem12Refutation.lean}.  The corrected
headline is \code{Theorem_1_2_intervalDomain_unconditional}.

\subsection{E4: Theorem 1.3(iv), exponent-domain gap}

The printed alternative (iv) assumes
\[
  \beta\ge\frac12,\qquad \alpha=2m+\gamma-2.
\]
Section~5.4 invokes Proposition~2.2 at
\[
  s(P)=\frac{P+\alpha}{\gamma},
\]
but that proposition requires \(s(P)>1\).  At the critical identity this is
equivalent to
\[
  P>2-2m.
\]
The proof starts from
\[
  q_*=\max\left\{1,\frac{N\alpha}{2}\right\},
\]
so its stated route needs the missing condition
\[
  \boxed{q_*>2-2m.}
\]

This condition is not implied by the printed hypotheses.  For example,
\[
  N=1,\quad m=\frac14,\quad\gamma=\frac72,\quad
  \alpha=2,\quad\beta=1
\]
satisfies \(\alpha=2m+\gamma-2\), while
\[
  q_*=1,\qquad s(q_*)=\frac67<1.
\]
Moreover \((N\alpha-2)_+=0\), so the first smallness disjunct in (1.25) is
automatic.  The printed argument therefore invokes Proposition~2.2 outside
its domain.  This is a proof gap, not a counterexample to boundedness.

The proof-supported amendment is
\[
  \max\{1,N\alpha/2\}>2-2m.
\]
The formalization records the corrected paper-constant route in
\code{Theorem_1_3_intervalDomainM_caseIV_with_paperConstants}.  It also closes
the printed theorem's actual one-dimensional conclusion for \(m<1\) by a
separate guard analysis: the global conjunct is printed only for \(m\ge1\),
while the bounded local branch is proved by
\code{Theorem_1_3_intervalDomainM_of_m_lt_one}.  The combined capstone is
\code{Theorem_1_3_intervalDomainM_caseIV_all_m}.  This alternative closure
does not make the displayed Section~5.4 invocation valid; it explains why
the repository can prove the final disjunct without silently endorsing that
step.

\subsection{E5: Proposition 1.1, maximal-continuation wording}

The finite endpoint alternative must be attached to the distinguished maximal
continuation.  Read instead as a dichotomy for every positive finite local
horizon \(T\), it is false: for \(a,b>0\), the positive logistic equilibrium
\[
  u(t,x)=\left(\frac ab\right)^{1/\alpha}
\]
is a classical solution on every finite interval, yet it neither blows up nor
approaches zero there.  The corrected statement constructs the supremum of
reachable horizons and asserts the blow-up/floor-loss alternative only when
that maximal horizon is finite; otherwise it glues the compatible local
solutions into a global branch.

This is a wording and carrier clarification.  The Lean witness against the
over-literal producer is
\code{not_legacyFiniteHorizonAlternativeProducer_of_positive_equilibrium};
the corrected endpoint-tail statement is
\code{correctedProposition_1_1_intervalDomainM} in
\code{Paper2/IntervalDomainMMaximalContinuationAlternative.lean}.

\subsection{E6: Lemma 2.7, missing control at the initial endpoint}

The literal statement assumes a differential inequality only for
\(t\in(0,\Tmax)\), but concludes a supremum bound on
\([0,\Tmax)\).  Take \(\Tmax=1\) and a spatially constant function
\[
  u(t,x)=t^{-1},\qquad t\in(0,1).
\]
With \(p=2\), its mass quantity is \(Y(t)=t^{-2}\), and
\[
  Y'(t)=-2t^{-3}\le -t^{-3}.
\]
Thus the interior differential inequality is satisfied for a suitable
choice of the lemma's nonnegative constants (the formal witness uses
\(C_1=C_2=C_3=0\), \(C_4=1\), and unit auxiliary exponents), while
\[
  \sup_{0<t<1}Y(t)=\infty.
\]
No interior differential inequality can control an omitted initial trace.

The corrected form assumes continuity, or merely a finite bound, for the
relevant \(L^p\)-mass at \(t=0\), then applies the scalar comparison on the
closed interval.  The exact certificates are
\code{not_Lemma_2_7_intervalDomain} and
\code{Lemma_2_7_intervalDomain_from_continuous_bound} in
\code{Paper2/IntervalDomainLemma27Unconditional.lean}.  In the PDE
application the classical initial trace supplies precisely this missing
endpoint information.

\section{Part II: persistence and stabilization}

\subsection{E7: Theorem 2.1(1), omitted pure-decay branch}

The printed part (1) assumes only \(m\ge1\) and claims a positive asymptotic
spatial lower bound for every global bounded positive solution.  The proof,
however, treats \(a=b=0\) and \(a>0,b>0\), omitting the allowed regime
\(a=0<b\).

For \(c_0>0\), define
\[
 u(t,x)=\bigl(c_0^{-\alpha}+\alpha bt\bigr)^{-1/\alpha},
 \qquad
 v(t,x)=\frac{\nu}{\mu}u(t,x)^\gamma.
\]
Again all spatial derivatives vanish, and the parabolic equation becomes
\[
  u_t=-bu^{1+\alpha}.
\]
The orbit is global, bounded, and strictly positive for every finite time,
but
\[
  \lim_{t\to\infty}\inf_{x\in\Om}u(t,x)=0.
\]
Thus the printed part (1) is false in the pure-decay regime.

The proof-faithful guard is
\[
  (a=b=0)\quad\text{or}\quad(a>0\land b>0).
\]
When \(a>0,b=0\), a globally bounded positive solution cannot exist because
its mass grows like \(e^{at}\), so the implication is vacuous.  The explicit
Lean refutation is
\code{not_Theorem_2_1_part1_intervalDomain_pureDecay}; the corrected statement
is \code{Theorem_2_1_part1_corrected}, assembled into
\code{Theorem_2_1_corrected_intervalDomainM}.

\subsection{E8: Theorems 2.2--2.5, all-time \(C^1\) decay}

Equation (2.12) asserts a \(C^1\) exponential estimate for every \(t\ge0\).
Theorem~2.2 assumes only \(u_0\in C(\overline\Om)\), positive and close in
\(L^\infty\).  Theorems~2.3--2.5 quantify over global bounded positive
solutions and reuse (2.12), again with no uniform \(C^1\) control at the
initial time.

On \(\Om=(0,1)\) with Neumann boundary conditions, take
\[
  u_{0,N}(x)=u_*+N^{-1/2}\cos(N\pi x).
\]
For large \(N\) the datum is positive, has the exact minimal-model mass, and
\[
 \norm{u_{0,N}-u_*}_{L^\infty}=N^{-1/2}\longrightarrow0,
 \qquad
 \norm{\partial_xu_{0,N}}_{L^\infty}=\pi N^{1/2}\longrightarrow\infty.
\]
Therefore an \(L^\infty\)-ball cannot have a uniform \(C^1\) bound at
\(t=0\).  For merely continuous \(u_0\), the \(C^1\) norm at \(t=0\) need not
even be defined.  Analytic-semigroup smoothing naturally has a
\(t^{-1/2}\)-type singularity as \(t\downarrow0\).

The correct statement separates two regimes:
\begin{enumerate}[leftmargin=2em]
\item if the initial perturbation is small in a fractional-power space
  \(X^\theta\hookrightarrow C^1\), then the sectorial stability theorem gives
  an all-time strong-norm estimate;
\item from \(L^\infty\)-small or merely continuous data, one first evolves to
  some \(t_0>0\), and then obtains \(C^1\) exponential convergence for
  \(t\ge t_0\).
\end{enumerate}
The linear stability/instability threshold itself is unaffected.

The interval-domain obstruction is
\code{not_intervalDomain_Theorem_2_5_original_allTime} in
\code{Paper3/IntervalDomainSectorialCorrectedObstruction.lean}.  The public
headline theorems in the repository use
\code{EventualExponentialC1ConvergenceWith} and orbit-dependent positive entry
times.  The same correction applies wherever Theorems~2.3--2.5 invoke
(2.12).

\section{Findings deliberately excluded from the paper errata}

Formal proof assistants make over-quantification easy to detect.  The
following refutations are valuable interface audits, but they must not be
reported as counterexamples to the source papers.

\begin{enumerate}[leftmargin=2em,itemsep=5pt]
\item \textbf{Traveling-wave Lemma 4.1.}
  The Lean proposition \code{Lemma_4_1} in
  \code{Paper1/Statements.lean} omitted the printed negative-branch speed
  threshold and upper bound on \(M\), and combined the two separately stated
  barriers into a global minimum.  Consequently \code{not_Lemma_4_1} refutes
  that stronger Lean proposition.  Likewise,
  \code{not_differentiableAt_upperBarrier_of_interface} observes that the
  minimum is not \(C^1\) at its interface, whereas the paper states
  \(e^{-\kappa x}\) only on one side and \(M\) separately.  Neither theorem is
  a paper-level refutation.

\item \textbf{Traveling-wave Lemma 2.1.}
  \code{not_forall_Lemma_2_1} chooses fabricated norms and semigroup data in
  an abstract structure.  The paper concerns the concrete heat semigroup, so
  this is a test of an overly universal API.

\item \textbf{Part I Theorem 1.3 \code{FrontierData}.}
  \code{not_IntervalDomainTheorem13FrontierData_of_real_regime} refutes an
  earlier structure containing arbitrary total functions and a global
  extension field.  Those fields are not printed hypotheses.  The genuine
  paper issue is only the exponent-domain gap E4.

\item \textbf{Part I Lemma 2.6.}
  An early abstract bootstrap interface stored a gradient frontier strong
  enough to encode the desired endpoint and did not expose the initial
  \(L^p\)-mass continuity.  The repository proves
  \code{Lemma_2_6_intervalDomain_from_continuous_lp_bounds}, but it does not
  claim a mathematical counterexample to the PDE lemma as used in the paper.

\item \textbf{Positive-time mass and zero-time signal storage.}
  Earlier Lean trajectory records constrained the solution only for \(t>0\)
  while allowing the stored slices \(u(0)\) or \(v(0)\) to be changed.
  Refutations obtained by re-anchoring those slices diagnose the data model.
  The paper starts from a genuine initial datum and trace.  The corrected Lean
  predicates use conserved positive-time mass or an explicit initial trace.
\end{enumerate}

\section{Reproducibility and interpretation}

The calculations above were checked against the archived paper text, not
against paraphrased Lean comments.  The repository's root build checks all
headline imports.  The intended interpretation of a Lean theorem name is:
\begin{itemize}[leftmargin=2em]
\item a theorem of the form \code{not_X} is evidence against the paper only
  after the hypotheses of \code{X} have been compared with the printed
  hypotheses;
\item a proof of a corrected headline establishes the corrected mathematical
  statement, but does not retroactively validate a defective displayed step;
\item an axiom-clean proof certifies deduction from the formal statement; it
  does not by itself certify that the formal statement has the intended
  semantics.
\end{itemize}

\noindent This distinction is why E2, E3, E6, and E7 are presented with
explicit mathematical counterexamples, E1 and E4 as proof gaps, E5 and E8 as
statement repairs, and the five interface findings in a separate section.

\end{document}
